% =====================================================================
%  The Spectral--Topological Decoupling Theorem in Self-Supervised
%  Physical World Models
%
%  Nature Machine Intelligence -- Letters (draft)
%  All numbers in Table 1 are mean +/- std over three training seeds,
%  generated by benchmark.py / benchmark_long.py on this machine
%  (PyTorch 2.10, CPU, deterministic per-step data). The F-long
%  (Path-B) row is the single completed seed of the long-horizon arm;
%  the second seed stalled without completing a step (see Methods).
% =====================================================================
\documentclass[11pt]{article}

\usepackage[margin=1in]{geometry}
\usepackage{amsmath,amssymb,amsthm}
\usepackage{booktabs}
\usepackage{graphicx}
\usepackage{xcolor}
\usepackage{microtype}

\newcommand{\FILL}[1]{\textcolor{red}{[#1]}}
\newcommand{\R}{\mathbb{R}}
\newcommand{\E}{\mathbb{E}}
\newcommand{\tr}{\operatorname{tr}}
\newcommand{\diag}{\operatorname{diag}}
\newcommand{\grad}{\nabla}
\newcommand{\sym}{\operatorname{sym}}

\newtheorem{proposition}{Proposition}
\newtheorem{lemma}{Lemma}
\newtheorem{theorem}{Theorem}

\title{\bf The Spectral--Topological Decoupling Theorem in\\
Self-Supervised Physical World Models}
\author{Sehaj Randhir Singh\textsuperscript{1,*}\\[2pt]
\small \textsuperscript{1}Independent researcher; partial affiliation with
NYU Tandon\\ School of Engineering, New York University, Brooklyn, NY, USA\\
\textsuperscript{*}Corresponding author: \FILL{email}}
\date{\today}

\begin{document}
\maketitle

% =====================================================================
\begin{abstract}
Joint-embedding predictive architectures (JEPAs) are increasingly given
Hamiltonian or thermodynamic structure to enforce physically faithful
latent dynamics. We show this hope fails at the level of the
optimization. On Lorenz-63 and the 4-D hyperchaotic Lorenz flow,
diffeomorphically lifted into 64-dimensional latent spaces, five
structurally distinct latent predictors are trained and audited under a
chaos-aware protocol. Every constraint class fails on a property disjoint from its target: the
unconstrained predictor explodes ($\lambda_1=+28.2$); the rigid
Hamiltonian freezes and yet posts the best overlap score; naive
metriplectic dissipation pins $\tr(R)/n$ to zero; a divergence gate pins
the spectral sum to within $1\%$ while $\lambda_1$ runs $2.5\times$ high;
and differentiable QR alignment shapes only a finite-horizon proxy
($-0.18$), not the asymptotic exponent ($+2.85$). A second system
inverts the decoupling and destroys the second positive exponent a
$\lambda_1$-only objective cannot name. At 2000 steps every alive arm's
exponent climbs to $\approx{+}8$--$+20$; we state the decoupling as a
theorem with a mandatory evaluation standard.
\end{abstract}

% =====================================================================
\section{Introduction}
\label{sec:intro}

A growing line of work trains world models whose latent dynamics are
restricted to physical form --- Hamiltonian \cite{greydanus2019hamiltonian,
cranmer2020lagrangian}, port-Hamiltonian \cite{van2020port}, or
metriplectic/GENERIC \cite{morrison1986metriplectic,grmela1997dynamics,
ottinger1997dynamics} structure --- on the premise that an inductive bias
that provably conserves energy or increases entropy should yield stable,
physically faithful rollouts. The same promise underpins learned world
models \cite{ha2018world,hafner2020mastering,bruce2024genie}: within
joint-embedding predictive architectures (JEPAs)
\cite{lecun2022pathway,assran2023self,bardes2024revisiting,garrido2024learning},
the same idea is applied to the latent predictor: constrain the operator
$\dot z = f_\theta(z)$ to a structure that a physicist recognizes, and
hope the representation space inherits the physics.

This Letter reports a systematic failure of that premise, obtained not by
increasing capacity or data but by isolating the \emph{structural} effects
of each constraint class on a single, fixed, fully reproducible benchmark.
Our object of study is a generalized metriplectic system
\begin{equation}
\label{eq:metriplectic}
\dot z \;=\; J(z)\,\grad H(z) \;-\; R(z)\,\grad S(z),
\end{equation}
with $J^\top=-J$ (reversible/Poisson part), $R=R^\top\succeq 0$
(dissipation), and scalar potentials $H$ (energy) and $S$ (dissipation
potential). We train five latent predictors --- unconstrained MLP,
rigid Hamiltonian, naive metriplectic, fixed (divergence-free-$J$)
metriplectic with an invariant dissipation gate, and spectrally aligned
metriplectic with a differentiable QR Lyapunov tracker --- to predict the
nonlinearly lifted Lorenz-63 flow in a 64-dimensional JEPA latent space,
and we evaluate them under a protocol designed for chaotic targets:
boundedness of the autoregressive rollout, the leading Lyapunov exponent,
the phase-space contraction rate, attractor overlap, and dissipation
engagement. To separate properties of the constraint hierarchy from
properties of any single dataset, the identical protocol is run on a
\emph{second}, topologically distinct system: the 4-D hyperchaotic
Lorenz flow \cite{qi2008hyperchaos}, which shares the same
volume-contraction rate as Lorenz-63 but carries two positive Lyapunov
exponents --- so that a $\lambda_1$-only spectral objective is blind to
the second unstable direction by construction.

Each arm fails, and each fails on a property disjoint from the one its
constraint nominally controls. We argue that these five failures are not
engineering deficiencies but the constructive instances of a decoupling
structure: aggregate thermodynamic identities constrain spectral
\emph{sums}; finite-horizon spectral statistics constrain windowed
\emph{tangent maps}; and set-level distributional losses constrain
\emph{pushforward measures}. None of these levels controls the others,
and no conjunction of them certifies the conjunction of properties a
physical world model must have: {bounded, dissipative, non-degenerate,
dynamically faithful}. We state this as a theorem whose lemmas are
elementary (each is a one-line linear-algebra or ergodic-theory fact) and
whose force is empirical: the failure atlas exhibits every class failing
simultaneously on two benchmark systems, including the most constrained
arm on both.

\paragraph{Related work.}
Spectral objectives have been studied extensively in self-supervised
learning as anti-collapse devices --- spectral contrastive learning
\cite{haochen2021spectral}, alignment--uniformity analyses
\cite{wang2020understanding}, whitening objectives
\cite{ermolov2021whitening}, Barlow Twins \cite{zbontar2021barlow}, VICReg
\cite{bardes2022vicreg}, feature-collapse studies of ViTs
\cite{caron2021emerging}, rank-based quality metrics
\cite{garrido2023rankme}, and dimensional-collapse analyses
\cite{jing2022dimensional} --- but these control the \emph{representation}
covariance, not the \emph{dynamics} operator; collapse prevention is a
different certification from physical fidelity, and our arm C shows the
two to be compatible (a full-rank, frozen, dynamically dead latent
field). On the dynamics side, structure-preserving networks --- neural
ODEs \cite{chen2018neural}, SymODEN \cite{zhong2020symoden}, dissipative
SymODEN \cite{zhong2021dissipative}, Hamiltonian and Lagrangian networks
\cite{greydanus2019hamiltonian,cranmer2020lagrangian,toth2020hamiltonian,lutter2019deep},
dissipative Hamiltonian networks \cite{sosanya2022dissipative},
metriplectic networks
\cite{hernandez2021metriplectic}, GENERIC-formatted networks
\cite{zhang2021generic}, port-Hamiltonian learning
\cite{van2020port}, and Lyapunov-stable neural ODEs
\cite{kolter2019lyapunov} --- typically verify their structure on the
\emph{training-time} loss or on short rollouts; the spectral fine
structure of the learned field is rarely audited. The relevant audit
tools exist: Benettin\textendash{}QR spectra have long been used to
characterize chaotic recurrent networks \cite{engelken2020lyapunov,vogt2022lyapunov},
and finite-time exponents are standard diagnostics in dynamical systems
\cite{haller2001distinguished}. What is missing is the negative
statement we supply: that \emph{constraining} the latent field through
any of these aggregate or windowed statistics does not \emph{certify}
the spectrum the field actually realizes.

The paper is deliberately deflationary, in deliberate contrast to the
law-discovery program \cite{brunton2016discovering,iten2020ai,karniadakis2021physics}:
we do not claim to recover thermodynamics in latent space, nor that any
architecture can. We claim the opposite: that constraint-based
certification of latent physical dynamics is, as currently practiced,
unsound, and we specify the evaluation standard that would make it sound.

% =====================================================================
\section{Results}
\label{sec:results}

\subsection{Benchmark construction}
\label{sec:bench}

We use the Lorenz-63 system \cite{lorenz1963deterministic},
$(\sigma,\rho,\beta)=(10,28,8/3)$, whose reference invariants are the
leading Lyapunov exponent $\lambda_1^{\rm ref}\approx 0.9056$ per time
unit (empirically $0.8840$ on our integrator, recovered by the same
Benettin machinery used for evaluation) and the phase-space divergence
$\grad\!\cdot F_L=-(\sigma+1+\beta)=-13.667$, which equals the sum of the
Lyapunov spectrum; the non-leading sum is $\kappa^{\rm ref}=\sum_{i\ge2}
\lambda_i^{\rm ref}=\grad\!\cdot F-\lambda_1^{\rm ref}=-14.55$ under the
runtime convention ($\lambda_2^{\rm ref}=0$, $\lambda_3^{\rm ref}
\approx-14.57$; the literature-spectrum sum is $-14.57$, a $0.1\%$
difference). A fixed nonlinear lift $\phi$ and an element-wise
nonlinear warp $W$ embed the attractor diffeomorphically into $\R^{64}$;
because the composite is a diffeomorphism, the latent reference spectrum
is that of Lorenz (plus 61 neutral directions), so a faithful latent
model must reproduce $\lambda_1\approx0.9$ and a contracting flow. A JEPA
encoder $g_\xi$ (MLP backbone, 256 hidden units, SIGReg
\cite{bardes2024sigreg} anti-collapse regularizer) maps observations to
$z\in\R^{64}$; the dynamics operator $f_\theta$ is trained with a latent
predictive (alignment) loss, a motion-matching penalty, and --- for the
physics arms --- the appropriate structure loss (Methods). Training is
fully deterministic per step (fixed seed, per-step trajectory draws);
every number below is mean$\pm$std over three seeds.

\subsection{The failure atlas}
\label{sec:atlas}

Table~\ref{tab:atlas} reports, per arm, the quantities that define the
failure: boundedness of the 200-step evaluation rollout ($\max\|z\|$,
with the fraction of seeds whose rollout left the encoder's training
support), the leading Lyapunov exponent $\lambda_1$ of the learned field,
the non-leading spectral sum (contraction) $\kappa=\sum_{i\ge2}\lambda_i$,
the dissipation engagement $\tr(R)/n$, the field-activity ratio
$m=\langle\|f_\theta\|\rangle/\langle\|f_{\rm ref}\|\rangle$ (a
flatlining detector), the Chamfer attractor-overlap score, and --- for
the spectral arms --- the training-time finite-horizon proxy
$\hat\lambda_1$. All quantities are defined in Methods.

\begin{table}[t]
\centering
\caption{\textbf{The five-arm failure atlas.} $\lambda_1$ is the leading
Lyapunov exponent of the learned latent field (reference $+0.884$);
$\kappa=\sum_{i\ge2}\lambda_i=\tr-\lambda_1$ is the non-leading spectral
sum (reference $-14.55$, i.e.\ $\grad\!\cdot F-\lambda_1^{\rm ref}$ under
the runtime convention); $\tr(R)/n$ is the mean dissipation-matrix trace
normalized by the latent dimension (metriplectic arms only); $m$ is the
field-activity ratio (flatlining detector, reference $1$); Chamfer is the
symmetric mean nearest-neighbor distance between the predicted and
encoded attractor clouds in latent space (lower is better, reference
$0$); $\max\|z\|$ is the maximum rollout norm over the 200-step
evaluation, with the fraction of seeds whose rollout left the training
support ($\|z\|\gtrsim10^2$ vs.\ in-support $\|z\|\approx10$);
$\hat\lambda_1$ is the training-time finite-horizon QR proxy.
Mean\,$\pm$\,std over three seeds. ${}^\dagger$Rollout left the training
support ($\|z\|\gtrsim10^2$ vs.\ in-support $\|z\|\approx10$): in all three seeds for B and F, in two of three for E. ${}^\ddagger$Not
reportable: the free MLP's Jacobian carries no volume law and its
evaluation rollout escapes (see text). ${}^\S$Single seed: the
long-horizon arm is reported on the one completed training run; a second
run stalled without completing a step (no progress in 45 min on
identical code and hardware --- a QR-backpropagation pathology, see
Methods). The value in parentheses is the top-8 QR readout of
$\lambda_1$ at encoded points.}
\label{tab:atlas}
\begingroup\footnotesize\setlength{\tabcolsep}{3pt}
\begin{tabular}{lccccccc}
\toprule
Model & $\max\|z\|$ & $\lambda_1$ & $\kappa$ & $\tr(R)/n$ & $m$ & Chamfer & $\hat\lambda_1$ \\
\midrule
B & $1.8{\times}10^6{}^{\dagger}$ & $+28.2\pm1.3$ & ${}^{\ddagger}$ & --- & $1821\pm1881$ & $3.4{\times}10^3$ & --- \\
C & $10.1\pm0.3$ & $+0.011\pm0.002$ & $+0.01\pm0.00$ & --- & $0.005\pm0.001$ & $3.85\pm0.05$ & --- \\
D & $35.5\pm24.5$ & $+1.50\pm0.86$ & $-1.4\pm1.2$ & $0.0000\pm0.0000$ & $0.88\pm0.13$ & $8.7\pm3.3$ & --- \\
E & $175\pm135{}^{\dagger}$ & $+2.22\pm1.80$ & $-16.0\pm1.9$ & $3.77\pm2.15$ & $1.72\pm0.98$ & $23.0\pm4.8$ & --- \\
F & $288\pm288{}^{\dagger}$ & $+2.85\pm1.89$ & $-16.1\pm2.0$ & $4.68\pm3.13$ & $2.17\pm1.46$ & $25.8\pm6.6$ & $-0.183\pm0.001$ \\
F-long & $54.2{}^{\S}$ & $+0.18$ ($+0.007$)& $-13.8$ & $1.58$ & $0.63$ & $15.2$ & $-0.146$ \\
\midrule
reference & --- & $+0.884$ & $-14.55$ & --- & $1$ & $0$ & --- \\
\bottomrule
\end{tabular}
\endgroup
\end{table}

\begin{figure}[t]
\centering
\includegraphics[width=\textwidth]{fig1_atlas.png}
\caption{\textbf{The failure atlas.} Mean$\pm$std over three seeds per arm:
(a) leading Lyapunov exponent $\lambda_1$ of the learned field (dashed:
reference $+0.884$) --- only the frozen arm C is near it, for the wrong
reason; (b) non-leading spectral sum $\kappa=\tr-\lambda_1$ (dashed:
reference $-14.55$) --- E and F pin the sum while $\lambda_1$ runs free;
(c) dissipation engagement $\tr(R)/n$ --- exactly zero for D; (d) field-
activity ratio $m$ (dashed: reference 1, log scale) --- C flatlines;
(e) Chamfer attractor overlap (log scale) --- C posts the best score
while frozen (overlap trap); (f) training proxy $\hat\lambda_1$ (red
\textasciicircum, spectral arms only) vs evaluation $\lambda_1$ (black
$\bullet$) --- for F the proxy is negative while the evaluation exponent
is positive, in every seed.}
\label{fig:atlas}
\end{figure}

\paragraph{Model B --- unconstrained MLP: spectral blow-up.}
The free residual predictor is not volume-constrained in any way; its
Jacobian acquires a large positive leading exponent, $\lambda_1=+28.2\pm
1.3$ --- $32\times$ the reference rate --- and the rollout leaves the
encoder domain in all three seeds ($\max\|z\|$ of order $10^5$--$10^6$),
with Chamfer degrading to $3.4\times10^3$ and a field $1.8\times10^3$
times too active ($m=1821\pm1881$). No thermodynamic statement can be made
--- and none is enforced --- which is precisely the control condition the
physics arms are meant to beat.

\paragraph{Model C --- rigid Hamiltonian: trivial representation freeze
and the overlap trap.}
The Hamiltonian arm, structurally volume-preserving
($\grad\!\cdot(J\grad H)=\tr(J\,\grad^2H)=0$ for constant skew $J$),
cannot contract phase-space volume and therefore cannot represent a
dissipative attractor. The optimizer's response is not a wrong
Hamiltonian but \emph{no motion at all}: the learned field gradients
flatline ($m=0.005\pm0.001$, a field $\approx200\times$ weaker than the
reference), all exponents collapse to $\approx0$
($\lambda_1=+0.011\pm0.002$; the trace reads $+0.02$, exactly the
volume-preserving value), and the predictor posts the \emph{best} Chamfer
score of all arms ($3.85\pm0.05$) by mapping every latent state into the
same small cluster of the attractor cloud. The set-level Chamfer distance
rewards statistical overlap of point clouds and is blind to whether the
points are produced by dynamics or by a constant map
(Proposition~\ref{prop:overlap}). A frozen predictor is thus ``optimal''
under the metric most commonly used to claim attractor recovery.

\paragraph{Model D --- naive metriplectic: the $R=LL^\top$ dead saddle.}
With $R=LL^\top$ and $L$ zero-initialized, the loss is stationary in $L$
at $L=0$ ($\partial R/\partial L=0$ there), a dead saddle from which no
first-order optimizer escapes. The model ``survives'' --- $\lambda_1=
+1.50\pm0.86$ is bounded and positive, and the trace reads $+0.1\pm2.1$
(an essentially volume-preserving loop) --- by keeping dissipation
exactly off: $\tr(R)/n=0.0000$ to all printed digits in every seed.
Positive semidefiniteness of $R$ is guaranteed by construction;
\emph{nonzero} dissipation is not. The apparent closeness to the
reference exponent is itself a window artifact: at 2000 steps the
shadow exponent reads $+4.5$--$+19.7$ on the hyperchaotic system
(Methods), so the dead-saddle field is the most expansive in the
atlas, not the nearest to the reference.

\paragraph{Model E --- fixed metriplectic with divergence gate: spectral
decoupling.}
Closing the two loopholes of D --- fixing $J$ to a constant
divergence-free operator and initializing $L\neq0$ --- makes the
invariant dissipation gate bite: $\tr(R)/n=3.77\pm2.15$ (engaged) and the
trace--divergence proxy pins the spectral sum to $-13.80\pm0.17$,
within $1\%$ of the reference $-13.667$; the non-leading sum
$\kappa=-16.0\pm1.9$. But the gate constrains only the \emph{sum} of the
spectrum. The leading exponent, driven by the Rayleigh quotient of
dissipation along the least-stable direction, is uncontrolled:
$\lambda_1=+2.22\pm1.80$ --- $2.5\times$ the reference and \emph{larger}
than the naive arm's $1.50$ --- and the global metrics degrade with it:
$\max\|z\|=175\pm135$ (two of three seeds leave the support) and Chamfer
$23.0\pm4.8$, worse than the dead-saddle arm D despite dissipation being
on exactly as designed. Pinning a thermodynamic aggregate leaves the
microscopic spectrum free (Proposition~\ref{prop:trace}).

\paragraph{Model F --- spectrally aligned metriplectic, short horizon:
the finite-time-horizon trap.}
Arm F adds a differentiable QR Lyapunov tracker (Methods) with a
six-step training window ($0.06$ time units $\approx0.05$ Lyapunov
times). At the optimum the finite-time proxy is $\hat\lambda_1=
-0.183\pm0.001$ --- a contracting estimate, negative in every seed ---
while the 200-step evaluation exponent is $\lambda_1=+2.85\pm1.89$:
the windowed tangent map that training sees says ``stable'', the
asymptotic tangent map says ``chaotic at $3.2\times$ the reference
rate''. The QR spectral term is not even the dominant shaper of the
proxy: its loss \emph{rises} from $0.78$ to $1.13$ while the total
loss falls by $15\times$, because alignment and motion dominate --- yet
the sign flip is robust across seeds, and evaluating with the training
window itself would certify F as stable (Proposition~\ref{prop:finite}).
The divergence gate still holds ($\tr-\lambda_1=-16.1\pm2.0$; the
trace proxy is $-13.24\pm0.19$, within $4\%$ of reference), so the arm
satisfies the macroscopic identity while its leading exponent is
$3.2\times$ the reference and its rollout drifts farthest of the alive
arms ($\max\|z\|=288\pm288$, all three seeds outside the support;
Chamfer $25.8\pm6.6$). A finite-horizon Lyapunov statistic is a
functional of the windowed tangent map; it is not a bound on the
Oseledets spectrum.

\paragraph{Model F-long --- long-horizon spectral alignment: the trap
is not a short-window artifact.}
Lengthening the differentiable QR window tenfold (60 steps,
$0.6$ time units $\approx0.5$ Lyapunov times, truncated BPTT with the
frame detached every 10 steps) does not close the decoupling; it
inverts it. The 60-step proxy converges to $\hat\lambda_1=-0.146$
--- still far from the $+0.884$ target, with the spectral loss
\emph{rising} from $0.78$ to $1.04$ while the total loss falls
ninefold --- and the asymptotic exponent collapses to
$\lambda_1=+0.18$ (the evaluation QR readout gives $+0.007$): the
leading axis is overdamped into a near-degenerate, effectively
non-chaotic flow. Every macroscopic certificate now fires: the trace
gate is satisfied to $0.4\%$ ($-13.61$ vs.\ $-13.667$), dissipation
is engaged ($\tr(R)/n=1.58$), the rollout stays bounded
($\max\|z\|=54$), and the Chamfer score ($15.2$) beats both E and
F. The model is well-behaved and spectacularly wrong: it has learned
a contracting, non-chaotic flow for a chaotic world. Backpropagation
through the QR chain is the only mechanism that keeps the spectral
term from being fully ignored --- its pre-clip gradient norm reaches
$11.3$ mid-training, engaging the norm-1 clip at five of the seven
logged checkpoints (the non-spectral arms carry no such excursion);
the probe on the initialization returns
finite norms ($8.0{\times}10^{-3}$--$8.9{\times}10^{-3}$ across
full-unroll/TBPTT-10/per-step), showing the instability is a
trained-state phenomenon, not an initialization one; and a second
long-horizon run stalled without completing a single training step,
a further manifestation of the same fragility. The finite-time
horizon trap is thus quantitative: lengthening the window moves the
proxy, but no window-length functional of the training tangent map
--- and no combination with the trace gate --- controls the
asymptotic Oseledets exponent or the dynamical content of the flow.

\subsection{Generality: the same atlas on a second system}
\label{sec:lz4}

A single benchmark cannot separate properties of the constraint
hierarchy from properties of one dataset. We therefore repeat the
\emph{identical} protocol --- same five predictors, same optimizer
schedule, same evaluation --- on a second, topologically distinct
reference flow: the 4-D hyperchaotic Lorenz system \cite{qi2008hyperchaos}
\begin{equation}
\label{eq:lz4}
\dot x=a(y-x),\qquad\dot y=bx-y-xz+w,\qquad\dot z=xy-cz,\qquad\dot w=-dx,
\end{equation}
with $(a,b,c)=(10,28,8/3)$ and $d=2.0$. Its divergence is
$-(a+1+c)=-13.667$ --- \emph{identical} to Lorenz-63 --- but its
Oseledets spectrum carries two positive exponents: consistent-tangent
Benettin QR at $dt=0.002$ over 60 time units yields
$(\lambda_1,\lambda_2,\lambda_3,\lambda_4)=(+0.769,+0.098,\approx0,-14.59)$,
and the runtime shadow measurement at the evaluation $dt=0.01$ gives
$\lambda_1=+0.737$. The second positive exponent is the point: a
$\lambda_1$-only spectral objective cannot even \emph{name} the second
unstable direction, let alone control it. Table~\ref{tab:lz4} reports the
same columns as Table~\ref{tab:atlas} plus the QR-readout second
exponent $\lambda_2$.

\begin{table}[t]
\centering
\caption{\textbf{The failure atlas on the 4-D hyperchaotic Lorenz
system.} Same protocol and columns as Table~\ref{tab:atlas}, plus
$\lambda_2$ (the QR-readout second exponent; reference $+0.098$).
Reference leading exponent: runtime-measured $+0.737$; non-leading sum
$\kappa^{\rm ref}=\tr-\lambda_1=-14.40$ (same convention as
Table~\ref{tab:atlas}). Symbols as in Table~\ref{tab:atlas}:
${}^\dagger$ the only escaping arm on this system is B;
${}^\ddagger$ not reportable on an escaped rollout. Mean$\pm$std over
three seeds.}
\label{tab:lz4}
\begingroup\scriptsize\setlength{\tabcolsep}{2pt}
\begin{tabular}{lcccccccc}
\toprule
Model & $\max\|z\|$ & $\lambda_1$ & $\lambda_2$ & $\kappa$ & $\tr(R)/n$ & $m$ & Chamfer & $\hat\lambda_1$ \\
\midrule
B & $3.9{\times}10^6{}^{\dagger}$ & $+29.2\pm3.0$ & ${}^{\ddagger}$ & ${}^{\ddagger}$ & --- & $4699\pm4718$ & $8.4{\times}10^3$ & --- \\
C & $9.9\pm0.2$ & $+0.011\pm0.006$ & $+0.00\pm0.00$ & $+0.01\pm0.01$ & --- & $0.004\pm0.000$ & $3.09\pm0.12$ & --- \\
D & $57.4\pm22.5$ & $+2.22\pm0.93$ & $+0.27\pm0.01$ & $-1.7\pm2.2$ & $0.0000\pm0.0000$ & $0.74\pm0.07$ & $14.7\pm2.4$ & --- \\
E & $66.7\pm15.9$ & $+0.09\pm0.38$ & $-0.08\pm0.02$ & $-13.85\pm0.56$ & $1.77\pm0.34$ & $0.93\pm0.16$ & $15.2\pm0.5$ & --- \\
F & $62.2\pm9.0$ & $-0.01\pm0.20$ & $-0.08\pm0.01$ & $-13.31\pm0.31$ & $1.64\pm0.05$ & $0.88\pm0.08$ & $14.6\pm0.4$ & $-0.18\pm0.01$ \\
\midrule
reference & --- & $+0.737$ & $+0.098$ & $-14.40$ & --- & $1$ & $0$ & --- \\
\bottomrule
\end{tabular}
\endgroup
\end{table}

The atlas reproduces, system to system, and its headline numbers are
strikingly close. The unconstrained MLP blows up with the same
signature ($\lambda_1=+29.2\pm3.0$, $\max\|z\|$ of order $10^6$, the
only escaping arm on this system); the rigid Hamiltonian freezes into
the same trivial solution ($m=0.004$, $\lambda_1\approx0$) and again
posts the best Chamfer ($3.09$); the naive metriplectic arm survives
by pinning $\tr(R)/n$ to exactly zero; and the divergence-gated and
spectral arms again satisfy the macroscopic identity (contraction
$-13.8$ and $-13.3$ vs.\ the reference $-13.667$, within $1$--$3\%$)
while their leading exponents read $+0.09$ and $-0.01$ --- an order of
magnitude \emph{below} the reference $+0.737$, where on Lorenz the same
gate let them run $2.5$--$3.2\times$ \emph{above}. The decoupling
direction inverts with the system: the spectral sum is pinned, and the
leading exponent wanders where it will --- high on Lorenz, dead on
LZ4 --- always off the reference.

The spectrally aligned arm's proxy is again wrong, and the
convergence check shows the two systems are one phenomenon. On
Lorenz-63, $\hat\lambda_1=-0.183$ while the asymptotic exponent is
$+2.85$ (a sign flip). On LZ4 the proxy $\hat\lambda_1=-0.18$ and the
200-step evaluation exponent $-0.01$ \emph{agree in sign} --- the
standard window reads a contracting, over-damped flow, far from the
$+0.737$ target --- but this is the same finite-window artifact: at
2000 steps the exponent of both gated arms climbs to $+7.8$--$+9.7$
and the rollout leaves the encoder domain (steps 304--402, Methods).
The spectral objective drives the leading exponent toward zero on
every window it can see, and the field obliges within that window; the
asymptotic instability is merely slower to manifest than on Lorenz-63.
The distinctive new fact is $\lambda_2$. The reference flow has a
second positive exponent $+0.098$; the QR readout of every constrained
arm is negative or zero at the evaluation window ($-0.08\pm0.02$ for
E, $-0.08\pm0.01$ for F, $0.00$ for C, and only $+0.27$ for the dead-
saddle arm D, whose exponents are accidental). A $\lambda_1$-only
spectral objective cannot even \emph{name} the second unstable
direction, let alone preserve it: the arms that fail on $\lambda_1$
also destroy $\lambda_2$, and the one arm with an engaged constraint
that might have acted on it (F) is the one whose objective is blind to
it. The spectral decoupling and finite-time traps are therefore not
Lorenz-63 artifacts: they follow from the constraint classes
(Propositions~\ref{prop:trace}--\ref{prop:finite}) and fire on any
flow in the class, with the error direction set by the system, not by
the constraint.

\begin{figure}[t]
\centering
\includegraphics[width=\textwidth]{fig3_lz4.png}
\caption{\textbf{The failure atlas on the 4-D hyperchaotic Lorenz
system.} Same panels as Fig.~\ref{fig:atlas} (mean$\pm$std over three
seeds; dashed: reference $\lambda_1=+0.737$). The hatched bars on the
left panel are the QR-readout second exponent $\lambda_2$ (dotted:
reference $+0.098$): the arms that fail on $\lambda_1$ also destroy
the second unstable direction --- which the $\lambda_1$-only spectral
target cannot even name. Note the decoupling \emph{direction} within
the standard 200-step window: constraining the trace pins the sum
($\kappa$ panel) while the leading exponent of the gated arms reads
$\approx0$ --- the mirror image of the Lorenz-63 atlas, where the same
constraints let $\lambda_1$ run $2.5$--$3.2\times$ too high. The
convergence check (Methods) shows both readings are finite-window
artifacts: at 2000 steps every alive arm's exponent climbs to
$\approx+8$--$+20$ and the gated arms escape, so the atlas's verdicts
are conservative.}
\label{fig:lz4}
\end{figure}

\subsection{The decoupling theorem}
\label{sec:theorem}

The five arms are the constructive content of a decoupling structure.
Let $f_\theta$ be a learned latent field and $A=Df_\theta$ its Jacobian.

\begin{proposition}[Aggregate gates pin only the mean; the dimension factor]
\label{prop:trace}
Let $\mathcal{R}_c=\{R\in\mathrm{Sym}(d):R\succeq0,\;\tr R=cd\}$ be the
feasible set of a trace (equivalently divergence, hence spectral-sum)
gate in $d$ latent dimensions. Then
\begin{equation}
  \sup_{R\in\mathcal{R}_c}\lambda_{\max}(R)=cd,
  \qquad
  \sup_{R\in\mathcal{R}_c}v^\top R\,v=cd\quad\text{for every unit }v,
\end{equation}
i.e.\ the gate pins the \emph{mean} eigenvalue $c$ but leaves the
Rayleigh quotient along any prescribed direction --- in particular along
the leading Oseledets direction $v_1$ --- free up to a factor $d$ above
the mean. Moreover, over the set of spectra with fixed sum $c$ (the
only thing a divergence gate constrains), the leading exponent is
\emph{unbounded above}: for every $M$ there is an admissible spectrum
with $\lambda_1>M$.
\end{proposition}

\begin{proof}
For PSD $R$, $\lambda_{\max}(R)\le\tr R=cd$ and $v^\top Rv\le cd$, and
both are attained by the rank-one matrix $cd\,uu^\top$ with $u=v$ (any
unit vector). For the spectrum statement, fix $c<0$ and $M>0$: the
tuple $(\lambda_1,\ldots,\lambda_d)=(M,0,\ldots,0,c-M)$ is admissible
(sum $c$, weakly ordered for $M\ge c$), with $\lambda_1=M$ arbitrarily
large; the negative tail $-M$ is absorbed by the most contracting
direction. The algebraic bound is the only one a sum constraint
provides; whether the optimizer \emph{realizes} the headroom is
empirical (arms E, F do: $\lambda_1=+2.2\ldots+2.85$ with the gate
satisfied). $\square$
\end{proof}

\begin{proposition}[Finite-time/asymptotic decoupling with rate]
\label{prop:finite}
Let $\Phi^t$ be a $C^{1+\alpha}$ flow with Oseledets top exponent
$\lambda_1$ and finite-time exponent
$\Lambda_T(x)=\frac{1}{T}\log\|D\Phi^T(x)\|$.
\begin{enumerate}
\item[(i)] \emph{Generic rate.} Under the standard fluctuation law for
the top exponent (positive variance $\sigma^2$ of the pointwise growth
rates),
\begin{equation}
  \Lambda_T(x)=\lambda_1+O\!\left(\sqrt{\tfrac{\log\log T}{T}}\right),
\end{equation}
so the finite-time error is $O(1)$ on sub-Lyapunov windows and remains
of order $10\%$ even at five Lyapunov times in the generic case. A
windowed functional of the tangent map cannot certify the asymptotic
exponent with useful accuracy at training horizons ($T\approx0.06$
Lyapunov times for the short arms).
\item[(ii)] \emph{Non-certifiability.} For every finite $T>0$ there
exist smooth flows with $\Lambda_T<0<\lambda_1$: a flow with an
arbitrarily long initial contraction epoch (a slow passage past a
stable region) followed by asymptotic expansion has negative
finite-time exponent over $[0,T]$ and positive Oseledets exponent. The
map $\theta\mapsto(\Lambda_T(\theta),\lambda_1(\theta))$ is therefore
not a graph over the parameter space, for any $T$; the differentiable
QR training signal $\Lambda_T(\theta)$ is not a certificate of
$\lambda_1(\theta)$.
\end{enumerate}
\end{proposition}

\begin{proof}
(i) is the log-law for top Lyapunov exponents
\cite{barreira2002lyapunov}: $T(\Lambda_T-\lambda_1)/\sqrt{2\sigma^2\log\log T}
\xrightarrow{d}\mathcal{N}(0,1)$ in the generic case. For (ii), let
$x_\varepsilon$ follow $\dot x=f(x)$ with a bounded region of
contraction of length $T$ preceding a region of expansion (existence by
smoothing a piecewise-linear flow); the finite-time exponent over
$[0,T]$ is negative while the time-mean growth rate from any generic
start is positive. $\square$
\end{proof}

\begin{proposition}[Set-metric/flow decoupling: the overlap trap]
\label{prop:overlap}
Chamfer-type distances are invariant under reordering of the point set.
For any target cloud $P\subset B(\mu,r)$, the constant map $f\equiv c$
with $c\in B(\mu,r)$ attains Chamfer distance $\le 2r$ while carrying
zero dynamical information.
\end{proposition}

\begin{proof}
Both nearest-neighbor terms collapse to radii about $c,\mu$. $\square$
\end{proof}

\begin{theorem}[Spectral--topological decoupling, no-go form]
\label{thm:no-go}
No constraint policy expressible as a combination of (i) aggregate
thermodynamic identities (trace/divergence gates), (ii) finite-horizon
spectral statistics (differentiable QR proxies), and (iii) set-level
distributional losses (Chamfer/overlap) certifies the conjunction
\{boundedness in the encoder domain, correct leading exponent,
non-degenerate field, dynamical fidelity\} for metriplectic latent
predictors of the form \eqref{eq:metriplectic}. The failure is
structural (Propositions~\ref{prop:trace}--\ref{prop:overlap}) and is
realized simultaneously by a five-arm atlas on \emph{two} benchmark
systems --- Lorenz-63 and the 4-D hyperchaotic Lorenz --- with each
constraint class failing on a property disjoint from the one it
enforces, including the fully constrained arm F on both systems.
\end{theorem}

The theorem is a statement about certification, not about existence: it
does not claim that no metriplectic latent model can match the
reference flows, only that the constraint classes currently used to
certify physical fidelity cannot, singly or jointly, guarantee the
properties they are invoked for. The propositions are elementary; their
content is that the three levels --- spectral sums, windowed tangent
maps, pushforward measures --- are mutually independent, and the atlas
exhibits all three traps firing on each of the two benchmark systems
(Fig.~\ref{fig:schematic}).

\begin{figure}[t]
\centering
\includegraphics[width=0.92\textwidth]{fig2_schematic.png}
\caption{\textbf{Spectral--topological decoupling.} Each constraint class
acts on one functional level --- (i) aggregate thermodynamic identities
(trace/divergence gates, spectral sums), (ii) finite-horizon spectral
statistics (differentiable QR, windowed tangent maps), (iii) set-level
distributional losses (Chamfer, pushforward measures) --- and each level
fails independently, producing three disjoint traps (arms E, F, C). No
combination of (i)--(iii) certifies the conjunction of properties a
physical world model must have (Theorem~\ref{thm:no-go}).}
\label{fig:schematic}
\end{figure}

% =====================================================================
\section{Discussion: two traps and a mandatory evaluation standard}
\label{sec:discussion}

\subsection{The overlap trap}
Set-level geometry is not dynamics. Chamfer and Hausdorff distances
between the predicted and encoded attractor clouds reward statistical
overlap; a predictor that has learned to output the attractor's centroid
region scores well while being dynamically dead. The literature's
practice of reporting attractor overlap without a mobility check is
therefore unsound: the frozen arm C posts the best Chamfer in our atlas
($3.85$) while carrying a field $\approx200\times$ weaker than the
reference. The corrective is not a better set metric but a gating
discipline: report overlap only conditioned on (i) a field-activity
floor, (ii) a non-vanishing leading exponent, and (iii) an ordered
(time-aligned) trajectory error --- or replace set metrics with
motion-invariant statistics.

\subsection{The finite-time-horizon trap}
A differentiable QR Lyapunov tracker is a legitimate optimization target
only if its horizon is matched to the physics it is meant to certify.
Over sub-Lyapunov horizons, finite-time exponents are dominated by local
stretch and rotation and are decoupled from the Oseledets flag; over a
Lyapunov-time horizon they become trainable but still constrain nothing
global. The trap is not a metric artifact: evaluating with the same
short window would certify the unstable arm F as stable --- its training
proxy is $-0.18$, the evaluation exponent $+2.85$, a sign reversal in
every seed. Lengthening the window tenfold inverts rather than closes
the trap (arm F-long): the proxy becomes consistent with the truth, but
the truth is now wrong --- the leading axis is overdamped to
$\lambda_1\approx0.2$ (QR readout $+0.007$), an effectively non-chaotic
flow that nonetheless satisfies the trace gate to $0.4\%$ and stays
bounded. Any spectral claim must therefore report the window-length
sensitivity, the proxy--target gap, and a long-horizon evaluation
spectrum.

\subsection{The through-line: no arm ever achieves the reference}
The atlas is stronger than the sum of its traps. \emph{No arm achieves
the reference leading exponent while remaining dissipative,
non-degenerate, and bounded --- on either system.} On Lorenz-63 the
unconstrained arm overshoots by $32\times$ and escapes; the
symplectic arm freezes; the naive
metriplectic arm's $+1.50$ at the evaluation window is the nearest to
$\lambda_1^{\rm ref}$, by keeping $R\equiv0$ --- yet its 2000-step
exponent is $+4.5$--$+19.7$, the most expansive field in the atlas;
the gated arm pins the sum and lets $\lambda_1$ run to $+2.2$; the
short-window spectral arm sign-flips its own proxy ($-0.18$ vs.\
$+2.85$); and the long-window spectral arm collapses the leading axis
to $\approx0$ within the evaluation window. On the hyperchaotic
system the same gate and the same spectral objective push the leading
exponent to $\approx0$ at the evaluation window from the
\emph{opposite} side --- again transiently: at 2000 steps both climb
past $+7.8$ and escape. Every lever a practitioner reaches for ---
more structure, a dissipation gate, a spectral objective, a longer
spectral window --- either leaves the leading exponent free or
destroys it, and the direction of the error is set by the system, not
the constraint, while the set-level and aggregate certificates that
are actually reported all fire successfully on at least one wrong arm.
This is the constructive content of Theorem~\ref{thm:no-go}: the
certificates and the property do not couple.

\subsection{A mandatory evaluation standard for physical world models}
\label{sec:standard}
For any claim that a learned latent dynamics is physically faithful, we
propose that the following be reported as a minimum:
\begin{enumerate}
\item the \emph{full} (or top-$k$, $k\ge 8$) Lyapunov spectrum over an
evaluation rollout of at least five Lyapunov times, not a single exponent
from a training-window estimate;
\item the \emph{proxy--target gap}: the training-horizon finite-time
exponent and the evaluation-horizon exponent, with the window-length
sensitivity;
\item \emph{domain invariance}: $\max\|z\|$ over the evaluation rollout
and the fraction of time spent outside the training support (escape
disqualifies regardless of losses);
\item a \emph{mobility guard}: the field-activity ratio with a collapse
floor (rejects overlap-trap freezes);
\item \emph{ordered} trajectory metrics alongside any set-level score;
\item a \emph{directional dissipation audit}: $\tr(R)/n$ and the Rayleigh
quotient $v_1^\top R v_1$ along the leading Oseledets direction --- the
spectral-decoupling trap is invisible to the trace alone;
\item \emph{atlas accountability}: an explicit statement of which trap
class the model is vulnerable to.
\end{enumerate}

\subsection{Limitations}
The theorem's force is the constructive atlas plus three elementary
lemmas; it does not rule out structure imposed \emph{inside} generative
decoders or data-driven degeneracy-aware brackets, which lie outside the
constraint hierarchy treated here. The residual violation of the GENERIC
degeneracy conditions ($J\grad S=0$, $R\grad H=0$) by the learned
potentials is itself a caveat on all metriplectic arms. The 61 transverse
reference exponents are neutral by construction, not by measurement.
All conclusions are for two dissipative chaotic flows of the same
family (Lorenz-type polynomial vector fields with identical divergence),
one encoder family (MLP), and one optimizer regime; the atlas is
evidence, and the propositions are the theorem. A reference flow from a
different dynamical class (e.g.\ a coupled-oscillator network with a
conserved quantity) would strengthen --- or, if the traps failed to
fire, weaken --- the claim, and is the immediate next experiment.

% =====================================================================
\section{Methods}
\label{sec:methods}

\subsection{Setting and reference dynamics}
\label{sec:methods-setup}

Let $M=\R^3$ carry the Lorenz-63 flow
$\dot x=F_L(x)$, $F_L=(\sigma(y-x),\,x(\rho-z)-y,\,xy-\beta z)$ with
$(\sigma,\rho,\beta)=(10,28,8/3)$. The reference invariants are the
leading Lyapunov exponent $\lambda_1^{\rm ref}\approx0.9056$ per time
unit and the constant divergence $\grad\!\cdot F_L=-(\sigma+1+\beta)=
-13.667$, which equals the sum of the full spectrum (the remaining
exponents are $\lambda_2=0$ and $\lambda_3\approx-14.57$). One Lyapunov
time is $T_L=\lambda_1^{-1}\approx1.13$ time units.

Observations are produced by a fixed nonlinear lift $\phi$ and an
element-wise nonlinear warp $W$,
\begin{equation}
  \phi(x)_j = m_{j0}\, \sin(\omega_j\, \tilde x_j) +
              m_{j1}\, \tanh(\tilde x_j^3) +
              m_{j2}\, \tanh(e^{\tilde x_j}-1),
  \qquad \tilde x = (x-\mu)/s,
\end{equation}
\begin{equation}
  W(u)=u+0.2\sin(2u)\cos(u^2),
\end{equation}
with fixed random weights $(m_j,\omega_j)$ and normalization chosen so
that each lifted feature lies in $[-1,1]$, on which $W$ has positive
derivative and is therefore a diffeomorphism of the bounded domain. The
composite $W\circ\phi$ embeds the attractor into $\R^{64}$; because the
lift is a diffeomorphism on the invariant manifold, the reference
spectrum of the latent flow is exactly that of Lorenz augmented by 61
neutral directions (Lyapunov spectra are invariant under smooth
conjugacies \cite{oseledets1968multiplicative,barreira2002lyapunov}). A
faithful latent model must therefore reproduce $\lambda_1\approx0.9$,
zero transverse exponents, and a volume-contracting flow with total
contraction $-13.667$.

For the second system, $M=\R^4$ carries the 4-D hyperchaotic Lorenz
flow \eqref{eq:lz4} with $(a,b,c,d)=(10,28,8/3,2.0)$, whose divergence
$-(a+1+c)=-13.667$ coincides with Lorenz-63 and whose Oseledets
spectrum $\lambda\approx(+0.769,+0.098,0,-14.59)$ (consistent-tangent
Benettin QR, $dt=0.002$, 60 time units, float64) carries two positive
exponents. The identical lift $\phi$ and warp $W$ map the 4-D
attractor into $\R^{64}$ (with fixed normalization and lift parameters
seeded separately), so the latent reference spectrum is the 4-D
spectrum augmented by 60 neutral directions. The F arm's spectral
target is the runtime-measured $\lambda_1=+0.737$ at the evaluation
$dt=0.01$, exactly as for Lorenz. Everything else --- encoder,
losses, curriculum, evaluation horizon, seeds --- is bit-identical
between the two systems; only the vector field, its normalization, and
the reference invariants differ. This is the sense in which the two
atlases test the constraint hierarchy rather than the dataset.

\subsection{Encoder, training, and regularization}
\label{sec:methods-encoder}

The encoder $g_\xi$ is an MLP backbone (vector mode, input dimension 32,
256 hidden units, GELU, LayerNorm) followed by a projector onto the
$d=64$ latent space. Anti-collapse is handled by SIGReg
\cite{bardes2024sigreg}: for $P=64$ random unit projections and $F=16$
frequencies $t$, the empirical characteristic function of the projected
batch is matched to the Gaussian characteristic function
$\exp(-t^2/2)$, an Epps--Pulley-style test that penalizes collapse to a
point or a subspace without any teacher or momentum encoder. The
prediction target is the latent $g_{\xi^-}(x_{t+\tau})$ from a
momentum-copied encoder; the dynamics operator is trained on the
multi-step alignment loss
\begin{equation}
  \mathcal L_{\rm align}=\E\big\| f_\theta^{\circ k}(z_t) - z_{t+k}
  \big\|^2, \qquad z_t=g_\xi(x_t),
\end{equation}
over a progressive-horizon curriculum $k\in\{1,3,6\}$ (30\%/60\% of the
training schedule). A motion-matching penalty with a detached target
magnitude prevents the trivial ``do-not-move'' solution that minimizes
alignment by freezing. Training: AdamW ($\mathrm{lr}=3\times10^{-4}$),
gradient clipping at norm 1, batch 128, 250 steps per seed, three seeds,
with fully deterministic per-step trajectory draws (per-step data seed),
so every reported number is bit-reproducible from (seed, schedule).

\subsection{The metriplectic tensor layout}
\label{sec:methods-metriplectic}

Physics arms parameterize the latent field on the generalized
metriplectic template \eqref{eq:metriplectic}: $J(z)=A(z)-A(z)^\top$
skew (reversible/Poisson part), $R(z)=L(z)L(z)^\top/n$ positive
semidefinite (dissipation), and $H(z),S(z)$ scalar potentials realized as
Softplus MLPs. Two structural facts anchor the analysis:
\begin{itemize}
\item For constant skew $J$, the reversible field is exactly
divergence-free, $\grad\!\cdot(J\grad H)=\tr(J\,\grad^2H)=0$ (skew
contracted with a symmetric Hessian), so all phase-space contraction
comes from the dissipative part and is controlled by $R$ alone.
\item The dissipation matrix is PSD by construction ($R=LL^\top$), but
its engagement is not: with $L$ zero-initialized the loss is stationary
in $L$ at $L=0$, a dead saddle from which no first-order optimizer
escapes and $R\equiv0$ forever (Arm D).
\end{itemize}
The strict GENERIC degeneracy conditions $J\grad S=0$ and $R\grad H=0$
\cite{grmela1997dynamics,ottinger1997dynamics} are deliberately not
imposed: the target is an open system that must lose mechanical energy,
and forcing $dH/dt=0$ would re-break dissipation. The residual
degeneracy violation is a documented caveat of every metriplectic arm.
The macroscopic control is the \emph{trace--divergence gate},
\begin{equation}
  \mathcal L_{\rm gate}=\beta_g\Big(\tfrac1n\tr R - r^\ast\Big)^2,
\end{equation}
which, via the divergence identity, pins the spectral \emph{sum}; it
cannot pin any individual exponent (Proposition~\ref{prop:trace}).

\subsection{The five arms}
\label{sec:methods-arms}

All arms share the encoder, optimizer, curriculum, and data; only the
structural constraint on $f_\theta$ differs.
\begin{itemize}
\item[\textbf{B}] \emph{Unconstrained MLP:} $f_\theta(z)=z+\mathrm{MLP}(z)$,
residual map, no structure, no physics loss (control arm).
\item[\textbf{C}] \emph{Rigid Hamiltonian:} $f_\theta(z)=J_c\grad H_c(z)$
with fixed constant skew $J_c$, quadratic $H_c$, and a flat
energy-conservation penalty (a closed, volume-preserving prior ---
structurally incapable of contraction).
\item[\textbf{D}] \emph{Naive metriplectic:} full
\eqref{eq:metriplectic} with learned $J$, $R=LL^\top$ zero-initialized,
no gate. The dead saddle keeps $\tr(R)/n$ at exactly zero.
\item[\textbf{E}] \emph{Fixed metriplectic:} constant divergence-free
$J$, nonzero $L$ initialization, trace--divergence gate. The gate bites;
the sum is pinned; $\lambda_1$ is not.
\item[\textbf{F}] \emph{Spectral-aligned (short horizon):} Arm E plus
the differentiable continuous-QR loss
$\beta_\lambda(\hat\lambda_1-\lambda_1^{\rm ref})^2$ with a six-step
window ($0.06$ time units $\approx0.05$ Lyapunov times).
\item[\textbf{F-long}] \emph{Spectral-aligned (long horizon):} Arm E
plus a 60-step differentiable QR loss ($0.6$ time units $\approx0.5$
Lyapunov times) with truncated BPTT (frame detached every 10 steps).
\end{itemize}
The physics arms use the dissipative energy-balance consistency penalty
(energy decay matched to dissipated power), and Arm C uses the flat
variant. The spectral weight is $0.10$ (F) and $0.20$ (F-long); the gate
weight is $0.02$.

\subsection{Differentiable continuous QR (Benettin)}
\label{sec:methods-qr}

Let $A(t)=\partial_z f_\theta(z(t))\in\R^{64\times64}$ be the tangent
operator and $X$ the fundamental matrix of $\delta\dot z=A(t)\delta z$,
factorized $X=QR$ with $Q\in O(64)$ and $R$ upper triangular,
$R_{ii}>0$. Substituting into $\dot X=AX$ and left-multiplying by
$Q^\top$ gives $Q^\top A Q = Q^\top\dot Q + \dot R R^{-1}$, a unique
decomposition into skew and upper-triangular parts, hence the
continuous-QR system
\begin{equation}
  \dot Q = Q\,\Omega,\quad \Omega_{ij}=
  \begin{cases}(Q^\top A Q)_{ij},& i>j\\0,& i=j\\
  -(Q^\top A Q)_{ji},& i<j\end{cases},\qquad
  \frac{d}{dt}\log R_{ii}=(Q^\top A Q)_{ii}.
\end{equation}
The diagonal of the pulled-back operator drives the exponents; the
strictly-lower part rotates the frame. The implementable discrete
scheme (Benettin--Galgani--Giorgilli--Strelcyn \cite{benettin1980lyapunov})
propagates an orthonormal frame $Q_k$, forms the Jacobian action
$J_{\rm local}Q_k$ by finite-difference JVP per column, advances
$V=Q_k+dt\,J_{\rm local}Q_k$, re-orthonormalizes $V=Q_{k+1}R_k$ by
Householder QR, and accumulates
\begin{equation}
  \hat\lambda_i(T)=\frac{1}{N\,dt}\sum_{k=1}^{N}\log\big|(R_k)_{ii}\big|,
  \qquad T=N\,dt.
\end{equation}
The chain is differentiable end to end (``torch.linalg.qr'' backprops
through the Householder reflections; the sign ambiguity of QR is fixed by
$\log|\cdot|$ with a floor). For the long-horizon arm the frame is
detached every $k=10$ steps (truncated BPTT). We probed backpropagation
through a 100-step QR chain under three regimes --- full unroll,
TBPTT-10, and per-step detach --- on the initialization used here; the
total gradient norms were $8.0\times10^{-3}$, $8.7\times10^{-3}$, and
$8.9\times10^{-3}$ respectively (all finite): on this initialization the
failure is not gradient explosion but the proxy--truth gap documented in
the Results. Two caveats are explicit: gradients are ill-conditioned
near $R_{ii}\to0$ or eigenvalue collisions, and --- decisively ---
$\hat\lambda_i(T)$ is a \emph{finite-time} exponent, not a bound on the
asymptotic Oseledets spectrum (Proposition~\ref{prop:finite}).

\subsection{Evaluation metrics}
\label{sec:methods-eval}

All evaluation is on held-out initial conditions (fixed seed), with a
200-step autoregressive rollout ($2$ time units $\approx1.8$ Lyapunov
times) in the learned latent space:
\begin{itemize}
\item \emph{Boundedness/domain invariance:} $\max\|z\|$ over the rollout;
escape from the training support ($\|z\|\gtrsim10^2$ vs.\ in-support
$\|z\|\approx10$) is scored per seed.
\item \emph{Leading Lyapunov exponent $\lambda_1$:} shadow-trajectory
Benettin on the one-step latent map, invariant under the coordinate
warp, directly comparable to $0.884$.
\item \emph{Contraction:} Hutchinson trace estimate
$\tr(J_{\rm map}-I)/dt=\E_\varepsilon[\varepsilon^\top(J_{\rm map}-I)\varepsilon]/dt$
(the divergence proxy, equal to the spectral sum) and its non-leading
part $\kappa=\tr-\lambda_1=\sum_{i\ge2}\lambda_i$.
\item \emph{Dissipation engagement:} $\tr(R)/n$ averaged over the
rollout (metriplectic arms).
\item \emph{Field-activity ratio $m$:} mean per-step displacement of the
predicted rollout divided by the (detached) target displacement --- a
flatlining detector; $m\approx1$ for a healthy field, $m\ll1$ for a
frozen one.
\item \emph{Attractor overlap:} symmetric Chamfer and Hausdorff
distances in latent space between the predicted rollout cloud and the
encoded ground-truth cloud.
\item \emph{Finite-horizon proxy:} the training-time $\hat\lambda_1$
read out at the end of training.
\end{itemize}
A full-spectrum (top-8) continuous-QR readout over the evaluation
rollout cross-checks $\lambda_1$ and the non-leading sum.

\subsection{Reproducibility}
\label{sec:methods-repro}

Every number in Tables~\ref{tab:atlas} and \ref{tab:lz4} is generated
by\texttt{benchmark.py} on CPU (PyTorch 2.10, Python 3.14): three
seeds, 250 training steps per arm, batch 128, deterministic per-step
data. The two systems are selected by \texttt{--system lorenz} and
\texttt{--system lz4}; the protocol and code path are otherwise
identical. Training is bit-reproducible for a fixed thread count; a
four-vs-eight-thread rerun of arm E shifted its exponent by
$|\Delta\lambda_1|\lesssim0.7$, an order of magnitude below the
seed-to-seed spread reported in the tables, so the conclusions are
insensitive to CPU float jitter. A fully independent retraining of
arms E and F on Lorenz-63 --- different wall-clock schedule and thread
count --- reproduced the atlas within seed-to-seed spread (E:
$\lambda_1=2.18\pm1.86$ vs.\ $2.22\pm1.80$; F: $2.93\pm1.85$ vs.\
$2.85\pm1.89$; $\tr(R)/n$ and Chamfer within $2\%$), and a third
independent retraining of the frozen arm C reproduced its statistics
(Chamfer $3.82$--$3.91$ vs.\ $3.85\pm0.05$; $\lambda_1\approx0.008$
vs.\ $0.011\pm0.002$). The convergence
check above reruns the shadow exponent on the saved weights, so the
tables, the convergence paragraph, and the reproducibility numbers are
drawn from independent artifacts of the same protocol.

\paragraph{Convergence of the reported exponents.}
The leading exponents in Tables~\ref{tab:atlas}--\ref{tab:lz4} are
200-step estimates (1.8 Lyapunov times). To certify that they are
asymptotic rather than transient, the shadow-trajectory exponent was
recomputed on the same final weights at 2000 steps (20 Lyapunov times)
for every checkpointed arm. The frozen arm C is the only converged
one: its exponents are $\approx0$ at both horizons on \emph{both}
systems ($|\lambda_1(2000)|\le0.018$ and $|\lambda_1(200)|\le0.018$,
all rollouts finite to step 2000). Every alive arm's 200-step value
understates its
2000-step value by an order of magnitude. On Lorenz-63, E grows from
$+0.8$--$+4.6$ to $+8.8$--$+9.4$ and F from $+1.7$--$+5.5$ to
$+9.8$--$+10.0$, with the latent trajectory leaving the encoder domain
at steps 213--300 of 2000 in all six seeds. On the hyperchaotic system
the pattern fires with a longer calm phase: the gated arms read
$\approx0$ at 200 steps but climb to $+7.8$--$+10.1$ and escape at
steps 304--407, and the dead-saddle arm D --- whose 200-step value is
the nearest to the reference --- reads $+4.5$--$+19.7$ at 2000 steps
(bounded trajectory, strongly expanding tangent map). The reported
200-step values therefore \emph{understate} the instability of every
alive arm on both systems, and the atlas's verdicts are conservative,
not optimistic: any evaluation window shorter than the instability's
onset time can certify a pathological field.

The full five-arm run takes approximately 6.2 wall-clock hours on a
20-core CPU under load. The long-horizon arm is
\texttt{benchmark\_long.py} with a 60-step horizon and truncated BPTT;
it is reported on the one completed seed (the five short-horizon arms
are three seeds each). A second long-horizon run stalled without
completing a training step (45 min without progress, CPU-bound, on
identical code and hardware), which we attribute to the
QR-backpropagation fragility documented in the Results --- an
additional, if anecdotal, instance of the same instability. The
empirical reference $\lambda_1=+0.8840$ is recovered by the same
Benettin machinery on the true Lorenz flow.

% =====================================================================
\section*{Data availability}
\begin{sloppypar}
All benchmark data (rollouts, checkpoints, evaluation JSONs, and
convergence artifacts for both systems) are generated by the released
code and are reproducible bit-for-bit from the stated seed schedules.
\end{sloppypar}

\section*{Code availability}
\begin{sloppypar}
The complete benchmark (five arms, evaluation protocol, and the
long-horizon variant) is available at
\texttt{https://github.com/sehajr-singhs/spectral-topological-decoupling}.
The five-arm atlas is reproduced by two commands: \texttt{python
benchmark.py --seeds 0 1 2 --steps 250} (Lorenz-63) and
\texttt{python benchmark.py --system lz4 --seeds 0 1 2 --steps
250} (4-D hyperchaotic Lorenz); the long-horizon variant by
\texttt{python benchmark\_long.py} with \texttt{--seeds 0 --steps
120 --horizon 60}.
\end{sloppypar}

\section*{Author contributions}
S.R.S.\ conceived the study, designed and implemented the benchmark
suite, ran all experiments, performed the analysis, and wrote the
manuscript.

\section*{Competing interests}
The author declares no competing interests.

\section*{Acknowledgments}
\FILL{}

\begin{thebibliography}{9}
\bibitem{lorenz1963deterministic} E.~N. Lorenz, Deterministic nonperiodic flow, \emph{J. Atmos. Sci.} \textbf{20}, 130 (1963).
\bibitem{grmela1997dynamics} M.~Grmela, H.~C. \"Ottinger, Dynamics and thermodynamics of complex fluids. I. Development of a general formalism, \emph{Phys. Rev. E} \textbf{56}, 6620 (1997).
\bibitem{ottinger1997dynamics} H.~C. \"Ottinger, M.~Grmela, Dynamics and thermodynamics of complex fluids. II. Illustrations of a general formalism, \emph{Phys. Rev. E} \textbf{56}, 6633 (1997).
\bibitem{morrison1986metriplectic} P.~J. Morrison, A paradigm for joined Hamiltonian and dissipative systems, \emph{Physica D} \textbf{18}, 410 (1986).
\bibitem{benettin1980lyapunov} G.~Benettin, L.~Galgani, A.~Giorgilli, J.-M.~Strelcyn, Lyapunov characteristic exponents for smooth dynamical systems and for Hamiltonian systems; a method for computing all of them, \emph{Meccanica} \textbf{15}, 9 (1980).
\bibitem{oseledets1968multiplicative} V.~I. Oseledets, A multiplicative ergodic theorem. Lyapunov characteristic numbers for dynamical systems, \emph{Trans. Moscow Math. Soc.} \textbf{19}, 197 (1968).
\bibitem{lecun2022pathway} Y.~LeCun, A path towards autonomous machine intelligence, preprint (2022).
\bibitem{assran2023self} M.~Assran, Q.~Duval, I.~Misra, P.~Bojanowski, P.~Vincent, M.~Rabinovich, Y.~LeCun, N.~Ballas, Self-supervised learning from images with a joint-embedding predictive architecture, \emph{CVPR} (2023).
\bibitem{bardes2024sigreg} A.~Bardes, J.~Ponce, Y.~LeCun, SIGReg: Sigmoidal Gaussian regularization for preventing collapse in joint-embedding predictive architectures, \emph{preprint} (2024).
\bibitem{greydanus2019hamiltonian} S.~Greydanus, M.~Dzamba, J.~Yosinski, Hamiltonian neural networks, \emph{NeurIPS} (2019).
\bibitem{cranmer2020lagrangian} M.~Cranmer, S.~Greydanus, S.~Hoyer, P.~Battaglia, D.~Spergel, S.~Ho, Lagrangian neural networks, \emph{ICLR} (2020).
\bibitem{van2020port} S.~A.~van~der~Schaft, Port-Hamiltonian systems, \emph{Encyclopedia of Systems and Control} (2020).
\bibitem{barreira2002lyapunov} L.~Barreira, Y.~Pesin, Lyapunov exponents and smooth ergodic theory, \emph{AMS Surveys and Monographs} \textbf{23} (2002).
\bibitem{haller2001distinguished} G.~Haller, Distinguished material surfaces and coherent structures in three-dimensional fluid flows, \emph{Physica D} \textbf{149}, 248 (2001).
\bibitem{qi2008hyperchaos} G.~Qi, G.~Chen, S.~Du, Z.~Chen, Z.~Yuan, Analysis of a new chaotic system, \emph{Physica A} \textbf{352}, 295 (2005). [4-D Lorenz-type hyperchaotic flow; the system used here is its standard (10, 28, 8/3, d) form.]
\bibitem{haochen2021spectral} J.~HaoChen, C.~Wei, A.~Gaidon, J.~Ma, Provable guarantees for self-supervised deep learning with spectral contrastive loss, \emph{NeurIPS} (2021).
\bibitem{zbontar2021barlow} J.~Zbontar, L.~Jing, I.~Misra, Y.~LeCun, S.~Deny, Barlow Twins: self-supervised learning via redundancy reduction, \emph{ICML} (2021).
\bibitem{bardes2022vicreg} A.~Bardes, J.~Ponce, Y.~LeCun, VICReg: variance-invariance-covariance regularization for self-supervised learning, \emph{ICLR} (2022).
\bibitem{jing2022dimensional} L.~Jing, P.~Vincent, Y.~LeCun, Y.~Tian, Understanding dimensional collapse in contrastive self-supervised learning, \emph{ICLR} (2022).
\bibitem{zhong2020symoden} Y.~D. Zhong, B.~Dey, A.~K. Chakraborty, Symplectic ODE-net: learning Hamiltonian dynamics with control, \emph{ICLR} (2020).
\bibitem{zhong2021dissipative} Y.~D. Zhong, B.~Dey, A.~K. Chakraborty, Dissipative SymODEN: encoding Hamiltonian dynamics with dissipation and control into deep learning, \emph{ICLR} (2021).
\bibitem{hernandez2021metriplectic} Q.~Hern\'andez, A.~Badias, D.~Gonzalez, F.~Chinesta, E.~Cueto, Deep learning of thermodynamics-aware reduced-order models from spatiotemporal data, \emph{Comput. Methods Appl. Mech. Eng.} \textbf{379}, 113749 (2021).
\bibitem{zhang2021generic} Z.~Zhang, S.~Shin, M.~Yacoub, J.~C.~Brigham, GENERIC formalism-based neural networks for learning dissipative and reversible dynamics, \emph{J. Comput. Phys.} \textbf{443}, 110520 (2021).
\bibitem{kolter2019lyapunov} J.~Z.~Kolter, G.~Manek, Learning stable deep dynamics models, \emph{NeurIPS} (2019).
\bibitem{engelken2020lyapunov} R.~Engelken, F.~Wolf, L.~F.~Abbott, Lyapunov spectra of chaotic recurrent neural networks, \emph{Phys. Rev. Lett.} \textbf{124}, 028301 (2020).
\bibitem{chen2018neural} R.~T.~Q. Chen, Y.~Rubanova, J.~Bettencourt, D.~Duvenaud, Neural ordinary differential equations, \emph{NeurIPS} (2018).
\bibitem{ha2018world} D.~Ha, J.~Schmidhuber, World models, \emph{arXiv:1803.10122} (2018).
\bibitem{hafner2020mastering} D.~Hafner, T.~Lillicrap, M.~Norouzi, J.~Ba, Mastering Atari with discrete world models, \emph{ICLR} (2021).
\bibitem{bruce2024genie} J.~Bruce, M.~Dennis, A.~Edwards, J.~Parker-Holder, Y.~Shi, E.~Hughes, M.~Mathet, A.~Beevi, F.~Hasselt, M.~Heess \emph{et al.}, Genie: generative interactive environments, \emph{ICML} (2024).
\bibitem{bardes2024revisiting} A.~Bardes, Q.~Garrido, J.~Ponce, X.~Chen, M.~Rabbat, Y.~LeCun, M.~Assran, N.~Ballas, Revisiting feature prediction for learning visual representations from video (V-JEPA), \emph{Trans. Mach. Learn. Res.} (2024).
\bibitem{garrido2024learning} Q.~Garrido, N.~Ballas, A.~Joulin, Y.~LeCun, Learning and leveraging world models in visual representation learning, \emph{arXiv:2403.00504} (2024).
\bibitem{wang2020understanding} T.~Wang, P.~Isola, Understanding contrastive representation learning through alignment and uniformity on the hypersphere, \emph{ICML} (2020).
\bibitem{ermolov2021whitening} A.~Ermolov, F.~Siomos, N.~Sebe, J.~Kannala, Whitening for self-supervised representation learning, \emph{ICML} (2021).
\bibitem{caron2021emerging} M.~Caron, H.~Touvron, I.~Misra, H.~J\'egou, J.~Mairal, P.~Bojanowski, A.~Joulin, Emerging properties in self-supervised vision transformers, \emph{ICCV} (2021).
\bibitem{garrido2023rankme} Q.~Garrido, R.~Balestriero, L.~Najman, Y.~LeCun, RankMe: assessing the downstream performance of pretrained self-supervised representations by their rank, \emph{ICML} (2023).
\bibitem{karniadakis2021physics} G.~E. Karniadakis, I.~G. Kevrekidis, L.~Lu, P.~Perdikaris, S.~Wang, L.~Yang, Physics-informed machine learning, \emph{Nat. Rev. Phys.} \textbf{3}, 422--440 (2021).
\bibitem{brunton2016discovering} S.~L. Brunton, J.~L. Proctor, J.~N. Kutz, Discovering governing equations from data by sparse identification of nonlinear dynamical systems, \emph{Proc. Natl Acad. Sci. USA} \textbf{113}, 3932--3937 (2016).
\bibitem{iten2020ai} R.~Iten, T.~Metger, H.~Wilming, L.~del Rio, R.~Renner, Discovering physical concepts with neural networks, \emph{Phys. Rev. Lett.} \textbf{124}, 010508 (2020).
\bibitem{toth2020hamiltonian} P.~Toth, D.~J. Rezende, A.~Jaegle, S.~Racani\`ere, A.~Botev, I.~Higgins, Hamiltonian generative networks, \emph{ICLR} (2020).
\bibitem{lutter2019deep} M.~Lutter, C.~Ritter, J.~Peters, Deep Lagrangian networks: using physics as model prior for deep learning, \emph{ICLR} (2019).
\bibitem{sosanya2022dissipative} A.~Sosanya, S.~Greydanus, Dissipative Hamiltonian neural networks: learning from chaotic trajectories with limited data, \emph{arXiv:2209.15832} (2022).
\bibitem{vogt2022lyapunov} R.~Vogt, M.~Puelma Touzel, E.~Shlizerman, R.~Engelken, On Lyapunov exponents for RNNs: understanding information propagation using dynamical systems theory, \emph{ICML} (2022).
\end{thebibliography}

\end{document}
